% Gerado a partir de Unidade_1/Resolucoes_e_Pratica/Lista_Treino_Dificil.md
\chapter{Lista de Treino - Nível Difícil (15 Questões)}

\textbf{Disciplina:} Sinais e Sistemas (TEC 412)\\
\textbf{Tópicos:} Análise Abstrata, Propriedades Complexas e Pegadinhas de Prova.

\section{Tópico 1: Sinais Pares e Ímpares}\label{tuxf3pico-1-sinais-pares-e-uxedmpares}

\textbf{Questão 1:} Prove que o sinal \(x(t) = t \sin(t) + \cos^2(t)\) é puramente par.

\begin{quote}
\textbf{Resolução:}\\
Aplicando o teste de reversão \(t \to -t\):\\
\(x(-t) = (-t) \sin(-t) + (\cos(-t))^2\)\\
Sabendo que \(\sin(-t) = -\sin(t)\) e \(\cos(-t) = \cos(t)\):\\
\(x(-t) = (-t)(-\sin(t)) + (\cos(t))^2\)\\
\(x(-t) = t \sin(t) + \cos^2(t) = x(t)\).\\
O sinal é inteiramente par, logo sua componente ímpar é obrigatoriamente nula para todo \(t\).\\
\textbf{Gabarito: Sinal 100\% Par.}
\end{quote}

\textbf{Questão 2:} A parte par de um sinal real é \(x_p(t) = |t|\). Sabendo que o sinal original \(x(t) = 0\) estritamente para \(t < 0\), recupere o sinal completo \(x(t)\).

\begin{quote}
\textbf{Resolução:}\\
Sabemos que a parte par é \(x_p(t) = \frac{x(t) + x(-t)}{2} = |t|\).\\
Para \(t > 0\), o sinal revertido \(x(-t)\) cai no território negativo. Como \(x(t)\) é zero para tempos negativos, \(x(-t) = 0\) para todo \(t>0\).\\
Assim, para \(t > 0\): \(\frac{x(t) + 0}{2} = |t| \implies \frac{x(t)}{2} = t \implies x(t) = 2t\).\\
Como ele é \(0\) para \(t < 0\), o sinal completo recriado é:\\
\textbf{Gabarito: \(x(t) = 2t \cdot u(t)\).}
\end{quote}

\section{Tópico 2: Energia e Potência}\label{tuxf3pico-2-energia-e-potuxeancia}

\textbf{Questão 3:} Encontre a potência do sinal da onda quadrada \(x(t) = \sum_{k=-\infty}^{\infty} rect\left(\frac{t - 2k}{1}\right)\).

\begin{quote}
\textbf{Resolução:}\\
Trata-se de um trem de pulsos retangulares. A largura do pulso é 1, e eles estão espaçados a cada 2 unidades de tempo. Logo, o período é \(T_0 = 2\).\\
A potência de sinais periódicos se calcula em um único período:\\
\(P = \frac{1}{2} \int_{-1}^{1} |x(t)|^2 dt\).\\
Dentro do intervalo de um período completo (\(t \in [-1, 1]\)), a onda retangular ``liga'' (vale 1) na faixa \([-0.5, 0.5]\) e fica zero no resto.\\
\(P = \frac{1}{2} \int_{-0.5}^{0.5} (1)^2 dt = \frac{1}{2} \left[ \tau \right]_{-0.5}^{0.5} = \frac{1}{2} \cdot 1 = 0.5 \text{ Watts}\).\\
\textbf{Gabarito: Potência de 0.5 W (Energia é \(\infty\)).}
\end{quote}

\textbf{Questão 4:} O sinal \(x(t) = \cos(t) \cdot u(t)\) é de energia ou de potência?

\begin{quote}
\textbf{Resolução:}\\
Apesar de começar no tempo zero devido ao degrau e ser nulo à esquerda, ele oscila até o infinito positivo com amplitude constante.\\
Sua energia diverge (área infinta). Indo pra potência calculada como o limite até \(T \to \infty\):\\
\(P = \lim_{T \to \infty} \frac{1}{2T} \int_{0}^{T} \cos^2(t) dt\).\\
Resolvendo e dividindo pelos \(2T\) infinitos, o valor da potência se assenta em exatamente metade da potência de uma senoide bidirecional total.\\
\textbf{Gabarito: Sinal de Potência (\(P = 0.25 W\)).}
\end{quote}

\section{Tópico 3: Memória}\label{tuxf3pico-3-memuxf3ria}

\textbf{Questão 5:} O sistema \(y(t) = x(2-t)\) possui memória?

\begin{quote}
\textbf{Resolução:}\\
Testando o instante presente \(t=0\): \(y(0) = x(2-0) = x(2)\).\\
O sistema, posicionado no instante \(0\), requisitou o valor que ocorrerá no instante \(2\). Diferiu do instante de leitura, obriga armazenamento de estado (ou antecipação de buffer).\\
\textbf{Gabarito: Com Memória.}
\end{quote}

\textbf{Questão 6:} O filtro não-linear de mediana local \(y[n] = \text{mediana}(x[n], x[n-1], x[n-2])\) exige memória?

\begin{quote}
\textbf{Resolução:}\\
O ato mecânico de extrair a mediana requer ordenar 3 amostras, sendo uma presente (\(n\)) e duas passadas (\(n-1, n-2\)). A leitura de instantes passados decreta a existência de registradores de atraso.\\
\textbf{Gabarito: Com Memória.}
\end{quote}

\section{Tópico 4: Causalidade}\label{tuxf3pico-4-causalidade}

\textbf{Questão 7:} O filtro de suavização avançada \(y(t) = \int_{t-1}^{t+1} x(\tau) d\tau\) é causal?

\begin{quote}
\textbf{Resolução:}\\
Avalie o limite superior da integral. Para determinar a saída no instante hoje \(t\), eu preciso integrar a área do sinal até o limite de amanhã (\(t+1\)). Como usa o conhecimento futuro do sinal\ldots{}\\
\textbf{Gabarito: NÃO-Causal.}
\end{quote}

\textbf{Questão 8:} O subamostrador contínuo da forma \(y(t) = x(t) \cdot \delta(t-3)\) é causal?

\begin{quote}
\textbf{Resolução:}\\
Cuidado: o ``3'' está na resposta impulsiva/função do sistema, mas o sinal na entrada exigido é APENAS o próprio instante de tempo presente: \(x(t)\). Em nenhum momento a função pediu para ler um \(x\) diferente.\\
\textbf{Gabarito: Causal.}
\end{quote}

\section{Tópico 5: Linearidade}\label{tuxf3pico-5-linearidade}

\textbf{Questão 9:} O sistema \(y(t) = x^2(t) + x(t)\) é linear?

\begin{quote}
\textbf{Resolução:}\\
Falha grotescamente na homogeneidade.\\
Seja \(x(t) \to y(t)\).\\
Aplique um ganho \(A\): entrada \(A \cdot x(t)\). A saída gerada será \((Ax)^2 + Ax = A^2x^2 + Ax\).\\
Para ser linear, a saída deveria ser estritamente \(A \cdot y(t) = A(x^2 + x) = Ax^2 + Ax\).\\
Como \(A^2x^2 \neq Ax^2\) para \(A \neq 1\), a regra quebra.\\
\textbf{Gabarito: NÃO-Linear.}
\end{quote}

\textbf{Questão 10:} O filtro de Volterra quadrático de 1ª ordem \(y[n] = x[n] \cdot x[n-1]\) preserva a superposição?

\begin{quote}
\textbf{Resolução:}\\
Seja entrada \(x_3[n] = x_1[n] + x_2[n]\) (omitindo ganhos pra ficar fácil).\\
\(y_3[n] = (x_1[n] + x_2[n]) \cdot (x_1[n-1] + x_2[n-1])\)\\
Ao fazer a distributiva, surgem termos cruzados (cross-talk):\\
\(x_1[n]x_1[n-1] + x_2[n]x_2[n-1] + x_1[n]x_2[n-1] + x_2[n]x_1[n-1]\).\\
Os dois primeiros termos são \(y_1[n] + y_2[n]\). Mas a equação gerou lixo extra e não linear cruzado entre os dois sinais de entrada.\\
\textbf{Gabarito: NÃO-Linear.}
\end{quote}

\section{Tópico 6: Invariância no Tempo}\label{tuxf3pico-6-invariuxe2ncia-no-tempo}

\textbf{Questão 11:} O reversor puro de tempo (fita de cassete rodando ao contrário) \(y(t) = x(-t)\) altera o formato se aplicar um atraso?

\begin{quote}
\textbf{Resolução:}\\
1) Atrase apenas a entrada criando \(x_a(t) = x(t-t_0)\). A saída processando \(x_a\) emite \(x_a(-t) = x(-t - t_0)\).\\
2) Atrase todo o sistema original global: Vá na fórmula \(y(t) = x(-t)\) e substitua a letra ``t'' por ``\((t-t_0)\)''.\\
\(y_2(t) = x(-(t-t_0)) = x(-t + t_0)\).\\
Observando a álgebra de sinais: \(x(-t - t_0) \neq x(-t + t_0)\). As direções de translação viraram opostas.\\
\textbf{Gabarito: Variante no Tempo.}
\end{quote}

\textbf{Questão 12:} O integrador comprimido \(y(t) = \int_{-\infty}^{2t} x(\tau) d\tau\) é invariante?

\begin{quote}
\textbf{Resolução:}\\
Atrasando a entrada geramos um sinal deslocado \(\tau-t_0\) dentro da integral. Mas quando aplicamos o atraso global no tempo, substituímos o limite superior para \(2(t-t_0) = 2t - 2t_0\). Como as escalas de atraso ficam diferentes (uma atrasa de \(t_0\), o limite de integração escorrega \(2t_0\)), os resultados divergem.\\
\textbf{Gabarito: Variante no Tempo.}
\end{quote}

\section{Tópico 7: Estabilidade BIBO}\label{tuxf3pico-7-estabilidade-bibo}

\textbf{Questão 13:} O sistema derivador ideal \(y(t) = \frac{dx(t)}{dt}\) garante estabilidade?

\begin{quote}
\textbf{Resolução:}\\
Para estabilidade BIBO, qualquer sinal limitado deve gerar saída limitada.\\
Imagine uma entrada \(x(t) = u(t)\) (a famosa função degrau, que é eternamente limitada, valor 0 e 1).\\
A derivada do degrau unitário é a função impulso de Dirac: \(\delta(t)\). O impulso tende à amplitude de \(+\infty\) no exato instante \(t=0\).\\
Entrou sinal limitado \(\to\) Produziu infinito.\\
\textbf{Gabarito: INSTÁVEL.}
\end{quote}

\textbf{Questão 14:} O filtro integrador atenuado por exponencial (``Filtro Leaky'') descrito por \(y(t) = \int_{-\infty}^{t} x(\tau) e^{-(t-\tau)} d\tau\) é estável?

\begin{quote}
\textbf{Resolução:}\\
Diferente do integrador puro que varre infinitamente constante 1, a função de atenuação decresce à medida que mergulhamos no passado. Assumindo que \(|x(\tau)| \le B\):\\
\(|y(t)| \le B \int_{-\infty}^{t} e^{-(t-\tau)} d\tau\). Resolvendo a integral da parte exponencial atenuadora resulta em exatamente \(1\). A saída máxima fica aprisionada em \(|y(t)| \le B \cdot 1 = B\).\\
Como nunca cresce, \textbf{Gabarito: É Estável.}
\end{quote}

\section{Tópico 8: Mistura Total}\label{tuxf3pico-8-mistura-total}

\textbf{Questão 15:} Dada a equação diferencial \(y(t) = \frac{dx(t)}{dt} + t \cdot x(t)\). Classifique.

\begin{quote}
\textbf{Resolução e Gabarito:}

\begin{enumerate}
\def\labelenumi{\arabic{enumi}.}
\tightlist
\item
  Memória: Derivada requer saber instantes de tempo minúsculamente próximos (delta). \textbf{Com Memória}.
\item
  Causalidade: A derivada causal retroativa só lê o instante anterior e o presente. \textbf{Causal}.
\item
  Linearidade: A derivada da soma é a soma das derivadas, e o produto distribui. \textbf{Linear}.
\item
  Invariância no Tempo: O multiplicador genérico livre `\(t\)' muda as regras da física do sistema no relógio. \textbf{Variante}.
\end{enumerate}
\end{quote}
